\section{Discussion}
\label{sec:discussion}

\subsection{Fidelity to AVO}

\Cref{tab:fidelity} compares avocode with the paper, object by object. Where the paper is specific,
avocode follows it: a single lineage, versions as git commits with scores, a vector $\f$ with a
correctness gate, a knowledge base, and a supervisor for stall and unproductive cycles. Where the
paper is silent, avocode has to choose, and it chooses the stricter option each time.

\begin{table}[t]
\centering\small
\caption{avocode compared with the AVO paper.}
\label{tab:fidelity}
\begin{tabularx}{\linewidth}{@{}l X X@{}}
\toprule
Aspect & AVO paper & avocode \\
\midrule
Domain & CUDA attention kernels on B200 & any repository with \file{.avo/score} \\
$\f$ & vector of TFLOPS per config; 0 if incorrect & JSON line with optional vector; \cmd{null} if $\neg\pass$, refused before ranking \\
Acceptance & ``matches or improves'' the best; vector rule and tolerance not given & \cref{def:commit}: dominance, strict gain on one config, relative floor $\varphi$; ties refused \\
Anti-gaming & not discussed & re-score at commit; no measuring less; frozen direction; bench hash gates \\
$\Pt$ storage & git commit with its score & commit $+$ trailers $+$ git note $+$ \file{lineage/vNNN.md} \\
Failed attempts & inside the agent's context, then lost & attempt log and content-keyed dead ends, read by later sessions \\
$\K$ & CUDA/PTX docs, FA4 source & \file{knowledge/} with provenance $+$ the lineage itself, hybrid search \\
Agent & one internal agent, unnamed LLMs & pi, Claude Code or Codex via portable skills \\
Continuity & one conversation for 7 days & fresh process per turn; state in the repository \\
Supervisor & a supervisor agent; mechanism not given & deterministic stall/thrash detector; citing directive; no LLM call \\
Population & single lineage; branching is future work & single lineage $+$ \avo{fan} probes, one promoted \\
Scale of evidence & 7 days, 40 versions, $>$500 directions & 65 minutes, 7 versions, 9 turns \\
\bottomrule
\end{tabularx}
\end{table}

\subsection{Strict acceptance and its cost}

Requiring a strict gain (\cref{eq:gain}) keeps the lineage a record of progress, makes ``best $=$
newest'' correct, and stops noise from committing non-changes. It has two costs that the case
studies make visible.

First, a monotone lineage cannot take a worse step to reach a better region. The project's own blog
post compares this to digging for diamonds and stopping just short: a promising idea whose first
implementation regresses becomes a dead end. Run~2 shows that a fresh agent can reach the same idea
again and implement it better, but nothing in the harness asks it to. Non-monotone acceptance in the
style of simulated annealing, or an archive of promising-but-refused candidates, would address this
at the price of a more complex lineage.

Second, the floor $\varphi$ interacts with the ``one idea per candidate'' rule. A real gain smaller
than $\varphi$ cannot commit alone, so the agent bundles it with a partner. That makes a later
refusal harder to attribute. A rule that requires both precision and small steps is in tension
whenever the steps are smaller than the noise. Better measurement, for example repeated sampling of
the candidate and re-measuring the incumbent on the same machine, would shrink $\varphi$ and ease
the tension. avocode does neither today; it compares a fresh measurement with a stored one.

\subsection{A deterministic supervisor}

The paper's supervisor is drawn as an agent. avocode's is a counter, a signature matcher and a
template. This is a deliberate trade. A deterministic supervisor is cheap enough to run after every
turn, cannot hallucinate a citation, and is easy to test at its exact firing point. It leaves the
judgment of novelty to the main agent and gives it evidence instead of advice: the versions it
must not re-derive, the dead ends it must not retry, and the documents nobody has read. The cost is
that it cannot recognize a repeated \emph{idea} that fails in a new way, or a plateau of tiny
committed gains. Whether this is enough is open, because the supervisor never fired in a real run.

\subsection{Externalized memory as the substrate of continuity}

The largest structural change from the paper is that avocode's agent is stateless. The paper's agent
carries its history in context for seven days, and the paper does not say how it handles context
limits. In avocode, every turn starts blank, and continuity lives in four files and a git history.
This has three effects. Any agent can be the operator, including a different one on each turn. A
crash loses at most one turn. And everything the loop learned can be audited, searched and reused by
a later session or a human. The cost is that each turn spends tokens re-reading the past; the
cache-read volume in run~2 (5.6 million tokens over three turns) reflects this.

\subsection{Agent-agnosticism}

Because the enforcement lives in the CLI and the instructions live in skills, the same harness runs
under three agents with no per-agent logic in the commit rule. Each adapter is a short command
template plus an output parser. The design forbids any capability that exists only as a pi tool, so the pi
extension can only make things faster or more convenient, never different.

\subsection{Self-hosting}

avocode was built by a loop of the same shape. An outer loop, \emph{ralph}, re-invoked a coding agent
with a fixed prompt: orient from \file{PLAN.md} and the last line of \file{PROGRESS.jsonl}, run the
health check, pick one task, prove it with a transcript in \file{evidence/}, commit, hand off and
exit. The agent had no memory between runs. The core harness was built in one run of about three
hours, and the remaining slices and four bug fixes in a second run of about four and a half hours,
for a recorded cost of \$129 over 18 iterations. The difference from \avo{run} is the fitness signal:
ralph's gate is ``lint, type check and tests pass'' plus a human-readable transcript, while \avo{run}'s
gate is a number and \cref{def:commit}. Several design points of \avo{run} came from running ralph,
including how it handles interrupts: an orphaned agent session once kept editing the checkout after
its loop was stopped, and two agents committed to one tree.
